% Class discussion: Odd Lots, "The Perfect Storm That Clobbered Hedge Funds" (Bloomberg, Nov 11 2026, recorded Nov 10).
% Hosts Joe Weisenthal and Tracy Alloway; guests Nishant Kumar (hedge funds) and Michael Ball (macro), Bloomberg News.
% One discussion question per slide. Timestamps refer to the episode.
% Source text: the AI-generated transcript on the Omny episode page. It garbles names ("Ballyasney" = Balyasny,
% "0.72" = Point72); fund figures are as stated on the show and were not checked against another source.
% Everything is paraphrased except one short quote (slide 4 title).
% Appendix numbers are back-of-envelope: A1 assumes 4% annual volatility and A3 assumes 6% PM volatility
% (assumptions, not from the episode). A3 uses the Brownian first-passage formula; a daily simulation gives 13% / 38%.
% The 0.45 vs 0.05 holding-period numbers come from conditional_strategies_figures.py.
\documentclass[aspectratio=169,10pt]{beamer}
\usetheme{default}
\usecolortheme{default}

\usepackage{graphicx}
\usepackage{adjustbox}
\usepackage{amsmath}
\usepackage{amssymb}
\usepackage{booktabs}
\usepackage{appendixnumberbeamer}  % Exclude appendix from slide count

\setbeamertemplate{navigation symbols}{}
\setbeamertemplate{items}[circle]

% Footline: Author, Event (Year) on left; page/total on right
\setbeamertemplate{footline}{
  \hbox{%
    \begin{beamercolorbox}[wd=.5\paperwidth,ht=2.5ex,dp=1ex,left]{author in head/foot}%
      \hspace*{2ex}\insertshortauthor, \insertshortinstitute\ (\insertdate)
    \end{beamercolorbox}%
    \begin{beamercolorbox}[wd=.5\paperwidth,ht=2.5ex,dp=1ex,right]{date in head/foot}%
      \insertframenumber/\inserttotalframenumber\hspace*{2ex}
    \end{beamercolorbox}%
  }
}

% Highlight boxes: the discussion question
\setbeamercolor{box1}{bg=red!12,fg=black}

% Where in the episode, and who said it
\newcommand{\src}[1]{{\footnotesize\textcolor{gray}{Episode #1}}\par\vskip 6pt}
% Which part of the course it connects to
\newcommand{\course}[1]{{\footnotesize\textcolor{gray}{Course link: #1}}}
% The discussion box
\newenvironment{discuss}{\begin{beamercolorbox}[rounded=true,sep=8pt]{box1}\textbf{Discuss}\par\vskip 3pt}{\end{beamercolorbox}}

\title{When Everyone Lost Together}
\subtitle{Class discussion: Odd Lots, ``The Perfect Storm That Clobbered Hedge Funds'' (Nov 11, 2026)}
\author[A. Moreira]{Alan Moreira}
\institute[Data-Driven Investing]{NYU Stern --- Data-Driven Investing}
\date{Fall 2026}

\begin{document}

\begin{frame}
\titlepage
\end{frame}

%================================================================== CONTEXT
\begin{frame}{October 2026: almost every hedge fund lost at the same time}
\vskip 4pt
\begin{itemize}
\item The big multi-strategy ``pod shops'' (Millennium, Balyasny, Point72, ExodusPoint) all lost money in the same month
\vskip 7pt
\item Funds concentrated in AI stocks lost up to a third; large macro funds had some of their worst months ever
\vskip 7pt
\item Then, two days after the midterms, the biggest single-stock selloff since Covid
\end{itemize}
\vskip 12pt
\begin{beamercolorbox}[rounded=true,sep=8pt,center]{box1}
Strategies that were supposed to be different all failed together. Why?
\end{beamercolorbox}
\vfill
{\footnotesize\textcolor{gray}{Odd Lots (Bloomberg), hosts Joe Weisenthal and Tracy Alloway, with Nishant Kumar (hedge funds) and Michael Ball (macro strategy), Bloomberg News. Recorded November 10, 2026.}}
\end{frame}

%================================================================== 1
\begin{frame}[label=q1]{Is down 2.3\% a bad month?}
\src{5:20 \textbullet\ Nishant Kumar}
\begin{itemize}
\item Millennium lost \textcolor{red}{2.3\%} in October. It averages about 14\% a year
\vskip 6pt
\item In the same month, Citadel and D.E.\ Shaw made money
\end{itemize}
\vskip 12pt
\begin{discuss}
\begin{itemize}
\item What do you need to know before calling a month extreme?
\item Does Citadel's positive month tell you anything about skill?
\end{itemize}
\end{discuss}
\vfill
\course{Performance Evaluation; Hiring and Firing} \hfill \hyperlink{appendix_z}{\beamerbutton{A number}}
\end{frame}

%================================================================== 2
\begin{frame}[label=q2]{``They paid for the hedge. They didn't get the hedge.''}
\src{6:00 \textbullet\ Nishant Kumar; 50:06 \textbullet\ Tracy Alloway}
\begin{itemize}
\item Pensions and sovereign funds do not hold hedge funds to match the S\&P 500 in good years
\vskip 6pt
\item They hold them to soften the blow when markets fall. In October they did not
\end{itemize}
\vskip 12pt
\begin{discuss}
\begin{itemize}
\item Is alpha against the market enough to judge such a fund, or do you need its beta in down months?
\item What is the fund's alpha relative to the portfolio the pension already holds?
\end{itemize}
\end{discuss}
\vfill
\course{Performance Evaluation; Capital Allocation II (factor bets and alpha bets)}
\end{frame}

%================================================================== 3
\begin{frame}[label=q3]{Was everyone implicitly long AI?}
\src{1:49, 10:27, 17:17 \textbullet\ Joe Weisenthal}
\begin{itemize}
\item Could you have made money in the last two years without being long AI?
\vskip 6pt
\item Torsten Slok (Apollo): the market has split into an AI factor and an anti-AI factor
\vskip 6pt
\item One PM is long chips; another is short Treasuries, betting AI capex drives inflation. Two trades, one bet
\end{itemize}
\vskip 10pt
\begin{discuss}
\begin{itemize}
\item How would you build an AI factor from stock data? If a fund's AI loading explains its returns, was that alpha?
\item A risk manager sees only each pod's returns. How would they find the common bet?
\end{itemize}
\end{discuss}
\vfill
\course{Sorts; Multi-Factor Models; Decomposing a Portfolio; PCA (coming)}
\end{frame}

%================================================================== 4
\begin{frame}[label=q4]{Correlations are low, until they are not}
\src{13:00, 18:25 \textbullet\ Nishant Kumar}
\begin{itemize}
\item Goldman: correlations between multi-strategy funds sat around \textcolor{red}{0.2--0.3} for years
\vskip 6pt
\item Kumar: correlation is not always visible. It jumps in months like October
\end{itemize}
\vskip 12pt
\begin{discuss}
\begin{itemize}
\item What diversification did investors think they were buying, and what did they get in October?
\item If you estimated the covariance matrix in calm years, what would your optimizer have done?
\end{itemize}
\end{discuss}
\vfill
\course{Portfolio Mathematics; Capital Allocation III (estimation error)} \hfill \hyperlink{appendix_div}{\beamerbutton{A number}}
\end{frame}

%================================================================== 5
\begin{frame}[label=q5]{Crowding, or momentum?}
\src{2:35, 12:29 \textbullet\ Tracy Alloway, Nishant Kumar}
\begin{itemize}
\item Alloway: are multi-strategy funds just momentum in disguise?
\vskip 6pt
\item Kumar: not quite. Ten funds can each like the same liquid stock for their own reasons. That is crowding, not chasing past returns
\vskip 6pt
\item But if they all change their minds at once, the crash looks the same
\end{itemize}
\vskip 10pt
\begin{discuss}
\begin{itemize}
\item What in the data would separate the two stories?
\item Does it matter to the investor which one is true?
\end{itemize}
\end{discuss}
\vfill
\course{Momentum (crashes)}
\end{frame}

%================================================================== 6
\begin{frame}[label=q6]{Lose 5\%, lose half your capital}
\src{14:55 \textbullet\ Nishant Kumar}
\begin{itemize}
\item At many pod shops: lose \textcolor{red}{5\%} and your capital is cut in half. Lose 7.5\% and you are out
\vskip 6pt
\item The ruthlessness is the point: no single PM can sink the firm
\end{itemize}
\vskip 12pt
\begin{discuss}
\begin{itemize}
\item Is this a test of a PM's skill? If not, what is it doing?
\item What happens when every firm's rule fires in the same week?
\end{itemize}
\end{discuss}
\vfill
\course{Hiring and Firing; Conditional Strategies I} \hfill \hyperlink{appendix_stop}{\beamerbutton{A number}}
\end{frame}

%================================================================== 7
\begin{frame}[label=q7]{\$2.2 trillion of positions on \$300 billion of capital}
\src{13:00 \textbullet\ Nishant Kumar, citing PivotalPath}
\begin{itemize}
\item The ten largest firms run about \$2.2 trillion of positions, including borrowed money
\vskip 6pt
\item The capital behind them is \$300--350 billion: roughly \textcolor{red}{6--7 times} leverage
\end{itemize}
\vskip 12pt
\begin{discuss}
\begin{itemize}
\item For a long-short book, what does gross leverage tell you about risk, and what does it hide?
\item What makes a hedged book lose on its longs and its shorts at the same time?
\end{itemize}
\end{discuss}
\vfill
\course{Leverage and Shorting (coming)} \hfill \hyperlink{appendix_lev}{\beamerbutton{A number}}
\end{frame}

%================================================================== 8
\begin{frame}[label=q8]{Why give money back?}
\src{23:32 \textbullet\ Nishant Kumar}
\begin{itemize}
\item Millennium sought \$20 billion, saw over \$30 billion of demand, and took \$22 billion
\vskip 6pt
\item Citadel, D.E.\ Shaw and Point72 have returned \textcolor{red}{\$35--40 billion} to investors
\vskip 6pt
\item How much do five more equity PMs add at a firm that already has dozens?
\end{itemize}
\vskip 10pt
\begin{discuss}
\begin{itemize}
\item Why would a fund with excess demand turn money away?
\item What does that say about how alpha scales with fund size?
\end{itemize}
\end{discuss}
\vfill
\course{Why Should This Work?; Capital Allocation I} \hfill \hyperlink{appendix_div}{\beamerbutton{A number}}
\end{frame}

%================================================================== 9
\begin{frame}[label=q9]{Did volatility control make the selloff worse?}
\src{30:14 \textbullet\ Michael Ball}
Who sold on the Thursday after the election:
\begin{itemize}
\item Dealers who had sold puts, selling stock as the market fell to stay hedged
\item Leveraged ETFs, which must sell more as prices fall
\item Volatility-control funds, forced to cut as volatility rose
\item Trend-following CTAs
\end{itemize}
\vskip 8pt
\begin{discuss}
\begin{itemize}
\item Volatility-control funds run the strategy from Conditional Strategies I. If enough money runs it, who is on the other side?
\item Does that explain why the strategy pays, or does it threaten the premium?
\end{itemize}
\end{discuss}
\vfill
\course{Conditional Strategies I (why would risk be paid less?)}
\end{frame}

%================================================================== 10
\begin{frame}[label=q10]{Reacting to every move}
\src{19:44, 41:01 \textbullet\ Tracy Alloway, Michael Ball}
\begin{itemize}
\item Funds cut AI exposure before the midterms, bought back on Wednesday, and were hit on Thursday. On Monday the rebound squeezed the shorts
\vskip 6pt
\item Our holding-period table: reacting daily to the last five days of volatility drops the market's managed Sharpe ratio to \textcolor{red}{0.05}, against 0.45 rebalancing monthly
\end{itemize}
\vskip 10pt
\begin{discuss}
\begin{itemize}
\item How fast should a risk rule react?
\item What do trading costs add to the answer?
\end{itemize}
\end{discuss}
\vfill
\course{Conditional Strategies I (holding periods); Transaction Costs (coming)}
\end{frame}

%================================================================== APPENDIX
\appendix
\begin{frame}
\begin{center}
{\Large Appendix: numbers for the discussion}
\end{center}
\end{frame}

\begin{frame}[label=appendix_z]{A1: How extreme is down 2.3\%?}
\begin{itemize}
\item Average month: $14\% / 12 \approx 1.2\%$
\vskip 6pt
\item Suppose annual volatility is 4\% (an assumption; the episode does not give it). A month's volatility is $4\% / \sqrt{12} \approx 1.15\%$
\vskip 6pt
\item Distance from the average: $\dfrac{-2.3\% - 1.2\%}{1.15\%} \approx \textcolor{red}{-3}$ standard deviations
\vskip 6pt
\item If returns were normal: about once every 60 years. Fat tails make it more frequent, but still rare
\end{itemize}
\vskip 10pt
\begin{beamercolorbox}[rounded=true,sep=8pt,center]{box1}
One month says a lot about risk and almost nothing about alpha.
\end{beamercolorbox}
\vskip 5pt
\hfill \hyperlink{q1}{\beamerreturnbutton{Back}}
\end{frame}

\begin{frame}[label=appendix_div]{A2: What more pods buy you}
Equal weights on $N$ pods, each with volatility $\sigma$ and pairwise correlation $\rho$:
\[ \sigma_{p} \;=\; \sigma\,\sqrt{\rho + \frac{1-\rho}{N}} \]
\begin{center}
\begin{tabular}{lcc}
\toprule
 & $\rho = 0.25$ (calm years) & $\rho = 0.8$ (October) \\
\midrule
10 pods & 0.57 $\sigma$ & 0.91 $\sigma$ \\
Infinitely many pods & \textcolor{red}{0.50 $\sigma$} & \textcolor{red}{0.89 $\sigma$} \\
\bottomrule
\end{tabular}
\end{center}
\vskip 6pt
\begin{itemize}
\item Past about ten pods, more pods barely help. The floor is $\sqrt{\rho}$
\item The 0.8 is illustrative: the episode says correlations jumped, not by how much
\end{itemize}
\vskip 5pt
\hfill \hyperlink{q4}{\beamerreturnbutton{Back to correlations}} \quad \hyperlink{q8}{\beamerreturnbutton{Back to capacity}}
\end{frame}

\begin{frame}[label=appendix_stop]{A3: How often does the 5\% rule fire?}
A PM running 6\% annual volatility (an assumption). Chance of being down 5\% of starting capital at some point within a year:
\begin{center}
\begin{tabular}{lc}
\toprule
PM's true Sharpe ratio & Chance the rule fires within a year \\
\midrule
1 (skilled) & \textcolor{red}{14\%} \\
0 (no skill) & \textcolor{red}{40\%} \\
\bottomrule
\end{tabular}
\end{center}
\vskip 6pt
\begin{itemize}
\item The rule fires on a skilled PM about a third as often as on an unskilled one: a cheap filter, not a skill test
\item Telling a Sharpe of 1 from 0 with a $t$-stat of 2 takes about 4 years: $t \approx \text{SR}\sqrt{T}$
\end{itemize}
\vskip 4pt
{\footnotesize Computed for a random walk with drift: $P\!\left(\min_{t\le T} X_t \le -a\right) = \Phi\!\left(\tfrac{-a-\mu T}{\sigma\sqrt{T}}\right) + e^{-2\mu a/\sigma^{2}}\,\Phi\!\left(\tfrac{-a+\mu T}{\sigma\sqrt{T}}\right)$, with $a = 5\%$, $\sigma = 6\%$, $\mu = \text{SR}\times\sigma$.}
\vskip 4pt
\hfill \hyperlink{q6}{\beamerreturnbutton{Back}}
\end{frame}

\begin{frame}[label=appendix_lev]{A4: Leverage arithmetic}
\begin{itemize}
\item Positions over capital: $2.2 / 0.35 \approx 6.3$ \quad to \quad $2.2 / 0.30 \approx 7.3$
\vskip 6pt
\item A long-short book that is half longs, half shorts loses nothing when the whole market falls
\vskip 6pt
\item It loses when the two sides move apart. If the longs fall 1\% and the shorts rise 1\%, the loss is 1\% of all positions: \textcolor{red}{6--7\% of capital}
\vskip 6pt
\item That is the crowded case: everyone sells the same longs and buys back the same shorts at once
\end{itemize}
\vskip 5pt
\hfill \hyperlink{q7}{\beamerreturnbutton{Back}}
\end{frame}

\end{document}
